\section{Цель работы}
Приобретение практических навыков создания динамических библиотек и программ, использующих их функции.

\section{Задание}
Создать две динамические библиотеки с одинаковыми контрактами и разными реализациями. Программа №1 должна использовать одну библиотеку на основании информации, полученной при компоновке. Программа №2 должна загружать библиотеки во время исполнения через интерфейс операционной системы и переключать реализации командой \texttt{0}. Команды \texttt{1} и \texttt{2} вызывают соответствующую функцию с переданными аргументами. Провести анализ двух способов использования библиотек~\cite{task}.

\subsection{Вариант 13}
Первая функция --- №2, вычисление производной $\cos x$:
\begin{verbatim}
float Derivative(float A, float deltaX);
\end{verbatim}
Реализация 1 использует разность вперёд, реализация 2 --- центральную разность:
\[
D_1(A,h)=\frac{\cos(A+h)-\cos A}{h},\qquad
D_2(A,h)=\frac{\cos(A+h)-\cos(A-h)}{2h}.
\]
Вторая функция --- №7, площадь фигуры:
\begin{verbatim}
float Square(float A, float B);
\end{verbatim}
Реализация 1 вычисляет площадь прямоугольника $S_1=AB$, реализация 2 --- прямоугольного треугольника $S_2=AB/2$, где $A$ и $B$ приняты за катеты.

\subsection{Формат ввода и ограничения}
\begin{itemize}
\item \texttt{1 A deltaX} --- вычислить производную; угол задаётся в радианах, шаг не равен нулю.
\item \texttt{2 A B} --- вычислить площадь; стороны строго положительны.
\item \texttt{0} --- переключить обе функции на другую реализацию в программе №2. В программе №1 выводится диагностика.
\end{itemize}
Каждая команда занимает отдельную строку. Пустые строки пропускаются. Неверное число аргументов, неизвестные команды, нечисловые и неограниченные значения отклоняются. Конец ввода завершает программу. Вычисления производной внутри библиотеки выполняются в \texttt{double}, результат возвращается как \texttt{float}. Это уменьшает ошибки округления, но не устраняет их при очень малом шаге.
